\chapter{Grecques et P\&L de couverture}\label{ch:dv:greeks-and-the-hedging-pnl}

Un desk d'options vend des \omterm{def:dv:greeks-and-the-hedging-pnl:cashgamma}{straddles} à un mois, à la monnaie, sur une action
à une volatilité de 20 et couvre leur delta chaque soir. Pendant les quatre
semaines suivantes, l'action bouge plus que d'habitude~; sa volatilité réalisée sur le mois est de 25. Le
premier matin, le rapport de risque du desk affichait deux chiffres, un \omterm{def:dv:greeks-and-the-hedging-pnl:greeks}{vega} de
23\,000~USD par point de volatilité et un \omterm{def:dv:greeks-and-the-hedging-pnl:greeks}{theta} de 7\,600~USD par jour en sa faveur,
et une ligne de calcul~: si l'action réalise cinq points de plus que ce qui a été
vendu, le book perd environ cinq fois son \omterm{def:dv:greeks-and-the-hedging-pnl:greeks}{vega}, soit 115\,000~USD. En fin de
mois, la perte est proche de ce montant, et le desk peut dire, jour après jour, d'où
elle vient~: chaque jour, le gamma de la position a coûté davantage que ce que le
\omterm{def:dv:greeks-and-the-hedging-pnl:greeks}{theta} a rapporté. Ce chapitre rend cette comptabilité exacte. Il nomme les
sensibilités, démontre l'identité qui relie le gamma au \omterm{def:dv:greeks-and-the-hedging-pnl:greeks}{theta}, la transforme en
formule pour le P\&L d'une option couverte, et mesure ce qu'ajoute une couverture
quotidienne plutôt que continue.

\section{Les grecques et leurs formes}

Le delta et le gamma d'une option sont définis dans le livre~1, chapitre~26~:
les dérivées première et seconde de sa valeur par rapport au sous-jacent.
La couverture en delta consiste à détenir l'opposé du delta en sous-jacent.

\begin{definition}[Grecques]\label{def:dv:greeks-and-the-hedging-pnl:greeks}
Les \emph{grecques}\index{grecques} d'une position sont les dérivées partielles de sa
valeur $V(t,S,\sigma,r)$ par rapport aux paramètres du modèle de valorisation~: le delta
$\Delta=\partial_SV$ et gamma $\Gamma=\partial_{SS}V$~; le \emph{vega}\index{vega}
$\mathcal V=\partial_\sigma V$~; le \emph{theta}\index{theta} $\Theta=\partial_tV$, la
variation de la valeur en fonction du temps calendaire, toutes choses égales par ailleurs~; le \emph{rho}\index{rho}
$\mathrm{Rho}=\partial_rV$~; et les sensibilités croisées d'ordre deux
\emph{vanna}\index{vanna} $\mathrm{Vanna}=\partial_{S\sigma}V$, soit la variation du delta
avec la volatilité, et le \emph{volga}\index{volga} $\mathrm{Volga}=\partial_{\sigma\sigma}V$,
la variation du vega avec la volatilité.
\end{definition}

Pour un call dans le \omterm{def:dv:black-scholes-three-ways:model}{modèle de Black--Scholes}, avec $\varphi$ la densité normale,
\[
\Gamma=\frac{e^{-qT}\varphi(d_1)}{S\sigma\sqrt T},\quad
\mathcal V=Se^{-qT}\varphi(d_1)\sqrt T,\quad
\mathrm{Vanna}=-\frac{e^{-qT}\varphi(d_1)\,d_2}{\sigma},\quad
\mathrm{Volga}=\mathcal V\,\frac{d_1d_2}{\sigma},
\]
et le put a le même gamma et le même \omterm{def:dv:greeks-and-the-hedging-pnl:greeks}{vega}. Les desks expriment les \omterm{def:dv:greeks-and-the-hedging-pnl:greeks}{grecques} en montants~: le \omterm{def:dv:greeks-and-the-hedging-pnl:greeks}{vega} par
point de volatilité (0.01), le \omterm{def:dv:greeks-and-the-hedging-pnl:greeks}{theta} par jour calendaire, le \omterm{def:dv:greeks-and-the-hedging-pnl:greeks}{rho} par point de base, et le gamma
par la grandeur suivante.

\begin{definition}[Cash gamma et straddle]\label{def:dv:greeks-and-the-hedging-pnl:cashgamma}
Le \emph{cash gamma}\index{cash gamma} d'une position est $\tfrac12\Gamma S^2$~: le P\&L
de la position couverte en delta par unité de rendement au carré, puisque
$\tfrac12\Gamma\,dS^2=\tfrac12\Gamma S^2(dS/S)^2$. Les rapports de risque donnent aussi
$\Gamma S^2/100$, la variation du delta en cash $\Delta S$ pour un mouvement de 1~\%. Un
\emph{straddle}\index{straddle} est un call et un put de même strike et même échéance~:
à la monnaie, son delta est proche de zéro et son gamma et son \omterm{def:dv:greeks-and-the-hedging-pnl:greeks}{vega} valent le double
de ceux du call.
\end{definition}

\begin{omfigure}
\begin{tikzpicture}
\begin{groupplot}[group style={group size=2 by 1,horizontal sep=1.6cm},
  width=6.6cm,height=4.8cm,xlabel={cours spot},
  tick label style={font=\small},label style={font=\small},
  xmin=70,xmax=130,grid=major,grid style={gray!25},table/col sep=comma]
\nextgroupplot[ylabel={gamma},ymin=0,ymax=0.095,
  legend columns=3,legend style={font=\footnotesize,at={(1.15,-0.32)},anchor=north,draw=none,fill=none,/tikz/every even column/.append style={column sep=8pt}}]
\addplot[omThm,very thick] table[x=spot,y=g1]{figdata/derivatives/04-greeks-and-the-hedging-pnl/shapes.csv};
\addplot[omDef,thick] table[x=spot,y=g025]{figdata/derivatives/04-greeks-and-the-hedging-pnl/shapes.csv};
\addplot[omProp,thick] table[x=spot,y=g005]{figdata/derivatives/04-greeks-and-the-hedging-pnl/shapes.csv};
\legend{1 an,3 mois,2.5 semaines}
\nextgroupplot[ylabel={vega par point},ymin=0,ymax=0.42]
\addplot[omThm,very thick] table[x=spot,y=v1]{figdata/derivatives/04-greeks-and-the-hedging-pnl/shapes.csv};
\addplot[omDef,thick] table[x=spot,y=v025]{figdata/derivatives/04-greeks-and-the-hedging-pnl/shapes.csv};
\addplot[omProp,thick] table[x=spot,y=v005]{figdata/derivatives/04-greeks-and-the-hedging-pnl/shapes.csv};
\end{groupplot}
\end{tikzpicture}

\omcaption{Gamma (à gauche) et \omterm{def:dv:greeks-and-the-hedging-pnl:greeks}{vega} par point de volatilité (à droite) d'un call de strike
100 ($\sigma=20\%$, $r=0$) à trois échéances. Les options courtes concentrent le gamma
près du strike~; les options longues portent le \omterm{def:dv:greeks-and-the-hedging-pnl:greeks}{vega}. Les deux sont positives pour toute
option longue. Données~: le code du chapitre.}\label{fig:dv:greeks-and-the-hedging-pnl:shapes}
\end{omfigure}

\section{L'identité gamma--theta et le P\&L de couverture}

\begin{proposition}[Identité gamma--theta]\label{prop:dv:greeks-and-the-hedging-pnl:gt}
La valeur de tout produit dérivé européen dans le \omterm{def:dv:black-scholes-three-ways:model}{modèle de Black--Scholes} vérifie
\[
\Theta+\tfrac12\sigma^2S^2\Gamma=r(V-S\Delta)+qS\Delta .
\]
Avec des taux et des dividendes nuls, $\Theta=-\tfrac12\sigma^2S^2\Gamma$~: une option longue
perd dans le temps exactement ce que son gamma rapporte lorsque l'action
bouge d'un écart-type.
\end{proposition}
\begin{proof}
Réarranger l'\omterm{def:dv:black-scholes-three-ways:pde}{équation de Black--Scholes} du chapitre~3.
\end{proof}

Maintenant, on considère l'option évaluée et couverte à une volatilité $\sigma$, tandis que l'action
évolue comme elle le fera. Sur un court intervalle, la position (option, $-\Delta$ actions,
financement monétaire des deux) varie par un développement de Taylor dans lequel le terme en delta est
neutralisé~:
\[
d\Pi=\Theta\,dt+\tfrac12\Gamma\,dS^2-r(V-S\Delta)\,dt-qS\Delta\,dt
=\tfrac12\Gamma S^2\Bigl(\bigl(\tfrac{dS}{S}\bigr)^2-\sigma^2dt\Bigr),
\]
d'après l'identité. C'est toute la théorie d'un book de volatilité en une ligne.

\begin{definition}[P\&L de couverture]\label{def:dv:greeks-and-the-hedging-pnl:hpnl}
Le \emph{P\&L de couverture}\index{P\&L de couverture} d'une option couverte en delta est la variation
de valeur de l'option, prise avec sa couverture et son financement. Pour une option longue
couverte à la volatilité $\sigma$, il vaut, à l'ordre deux sur chaque intervalle,
le \omterm{def:dv:greeks-and-the-hedging-pnl:cashgamma}{cash gamma} multiplié par l'écart entre le rendement au carré et la variance
que supposait la couverture~:
\[
\pnl_{[0,T]}=\int_0^Te^{r(T-t)}\tfrac12\Gamma_tS_t^2\bigl(\sigma_{\mathrm{real},t}^2-\sigma^2\bigr)dt ,
\]
où $\sigma^2_{\mathrm{real},t}dt=(dS_t/S_t)^2$ est la
variance réalisée instantanée.
\end{definition}

La formule dit trois choses. Le P\&L se gagne au rythme du gamma en espèces, de sorte
que les mouvements des jours où le gamma est grand comptent davantage que ceux des jours
où il est faible. Elle est quadratique en les rendements et aveugle à leur signe. Enfin,
elle fait intervenir la variance réalisée, et non la volatilité réalisée, un fait que le
chapitre~14 transforme en produit.

\begin{proposition}[P\&L de couverture espéré]\label{prop:dv:greeks-and-the-hedging-pnl:expected}
Si l'action suit le \omterm{def:dv:black-scholes-three-ways:model}{modèle de Black--Scholes} de volatilité $\sigma_{\mathrm{real}}$
et que l'option est vendue à $\sigma$ et couverte en continu à $\sigma$, le \omterm{def:dv:greeks-and-the-hedging-pnl:hpnl}{P\&L de couverture}
espéré de l'option longue (avec la dérive à $r$) vaut
$V(\sigma_{\mathrm{real}})-V(\sigma)$ en valeur actuelle.
\end{proposition}
\begin{proof}
La valeur actuelle du P\&L vaut $\E\bigl[\int_0^Te^{-rt}\tfrac12\Gamma_t(\sigma)S_t^2
(\sigma_{\mathrm{real}}^2-\sigma^2)dt\bigr]$. La fonction $u=V(\sigma_{\mathrm{real}})-
V(\sigma)$ résout l'équation de valorisation à $\sigma_{\mathrm{real}}$ avec le
terme source $\tfrac12(\sigma_{\mathrm{real}}^2-\sigma^2)S^2\Gamma(\sigma)$ et une valeur terminale nulle~;
la formule de Feynman--Kac (livre~4, chapitre~4) donne le résultat.
\end{proof}

\begin{example}[Le mois du desk]\label{ex:dv:greeks-and-the-hedging-pnl:month}
Le \omterm{def:dv:greeks-and-the-hedging-pnl:cashgamma}{straddle} à la monnaie à un mois sur une action cotant 100 vaut 4.6059 à une volatilité
de 20 et 5.7570 à 25, avec des taux nuls. Un desk vendeur de 1\,000 \omterm{def:dv:greeks-and-the-hedging-pnl:cashgamma}{straddles}
(multiplicateur 100) s'attend à perdre $(5.7570-4.6059)\times100\,000$, soit environ
115\,100~USD, lorsque la volatilité réalisée atteint 25. Son \omterm{def:dv:greeks-and-the-hedging-pnl:greeks}{vega}, 0.2302 par point
et par \omterm{def:dv:greeks-and-the-hedging-pnl:cashgamma}{straddle}, multiplié par cinq points, donne presque le même montant~: pour de petites variations
la perte espérée vaut le \omterm{def:dv:greeks-and-the-hedging-pnl:greeks}{vega} multiplié par l'écart de volatilité.
\end{example}

\begin{omfigure}
\begin{tikzpicture}
\begin{axis}[width=10.5cm,height=5cm,
  xlabel={jours de négociation},ylabel={P\&L par straddle},
  tick label style={font=\small},label style={font=\small},
  xmin=0,xmax=21,ymin=-2,ymax=0.4,grid=major,grid style={gray!25},
  legend columns=2,legend style={font=\footnotesize,at={(0.5,-0.3)},anchor=north,draw=none,fill=none,/tikz/every even column/.append style={column sep=8pt}},
  table/col sep=comma]
\addplot[omDef,very thick] table[x=day,y=pnl]{figdata/derivatives/04-greeks-and-the-hedging-pnl/gamma_path.csv};
\addplot[omThm,thick,dashed] table[x=day,y=pred]{figdata/derivatives/04-greeks-and-the-hedging-pnl/gamma_path.csv};
\legend{straddle vendu à découvert couvert,$-\sum\frac12\Gamma S^2\bigl(r_k^2-\sigma^2\Delta t\bigr)$}
\end{axis}
\end{tikzpicture}

\omcaption{Un mois d'un \omterm{def:dv:greeks-and-the-hedging-pnl:cashgamma}{straddle} vendu à découvert à une volatilité de 20, couvert quatre fois
par jour, alors que l'action réalise 25. La ligne en pointillés ajoute les termes gamma--thêta
intervalle par intervalle~; elle ne s'écarte jamais de plus de 0.09 du P\&L du livre
et se termine à 0.04 de celui-ci.
Le P\&L baisse par paliers les jours de forts mouvements et remonte les jours calmes,
lorsque le thêta l'emporte. Données~: le tutoriel.}\label{fig:dv:greeks-and-the-hedging-pnl:path}
\end{omfigure}

\section{Erreur de couverture discrète}

L'identité n'est exacte que dans la limite d'une couverture continue. Un desk qui couvre
une fois par jour détient un delta périmé entre deux rééquilibrages, et le rendement au carré
sur une journée est une estimation bruitée de la variance du jour.

\begin{definition}[Erreur de couverture discrète]\label{def:dv:greeks-and-the-hedging-pnl:dhe}
L'\emph{erreur de couverture discrète}\index{erreur de couverture discrète} est l'écart entre la
valeur terminale d'une option couverte à $N$ dates et celle de la même
option couverte en continu, lorsque la volatilité réalisée est égale à la volatilité de couverture.
Sa moyenne est proche de zéro et son écart type décroît en $1/\sqrt N$.
\end{definition}

\begin{proposition}[Ordre de grandeur de l'erreur]\label{prop:dv:greeks-and-the-hedging-pnl:dk}
Pour une option proche de la monnaie couverte $N$ fois à sa volatilité implicite $\sigma$,
que l'action réalise, l'écart type de l'erreur de couverture vaut
approximativement $\sqrt{\pi/4}\,\mathcal V\sigma/\sqrt N$.
\end{proposition}
\begin{proof}[Preuve partielle]
Sur un pas, l'erreur vaut $\tfrac12\Gamma S^2(x^2-\sigma^2\Delta t)$, où $x$ est le
rendement du pas~; sa variance vaut $\tfrac12\Gamma^2S^4\sigma^4\Delta t^2$ pour des rendements
gaussiens. En sommant les pas indépendants, avec $\Delta t=T/N$, et en moyennant
$\Gamma^2S^4$ sur les trajectoires, on obtient la décroissance en $1/\sqrt N$ et la dépendance en \omterm{def:dv:greeks-and-the-hedging-pnl:greeks}{vega}~; la
constante $\sqrt{\pi/4}$ provient de cette moyenne pour une option à la monnaie
(Derman et Kamal).
\end{proof}

Pour le \omterm{def:dv:greeks-and-the-hedging-pnl:cashgamma}{straddle} à un mois couvert chaque jour, la formule donne 0.890 et une simulation
de 20\,000 trajectoires donne 0.865~; une couverture huit fois par jour ramène la valeur simulée
à 0.308 (la \cref{fig:dv:greeks-and-the-hedging-pnl:discrete}). Le bruit est
aussi grand que l'avantage~: un desk qui vend la volatilité de deux points au-dessus de ce que
l'action réalisera gagne en moyenne deux fois le \omterm{def:dv:greeks-and-the-hedging-pnl:greeks}{vega} par \omterm{def:dv:greeks-and-the-hedging-pnl:cashgamma}{straddle}, soit environ 0.46, et
avec une couverture quotidienne, présente un écart type de 0.87 autour de cette valeur. Vendre de la volatilité
est un métier statistique, exercé sur de nombreuses positions.

\begin{omfigure}
\begin{tikzpicture}
\begin{axis}[width=10.5cm,height=4.8cm,
  xlabel={P\&L à l'échéance par straddle},ylabel={densité},
  tick label style={font=\small},label style={font=\small},
  xmin=-4,xmax=4,ymin=0,ymax=1.7,grid=major,grid style={gray!25},
  legend columns=2,legend style={font=\footnotesize,at={(0.5,-0.3)},anchor=north,draw=none,fill=none,/tikz/every even column/.append style={column sep=8pt}},
  table/col sep=comma]
\addplot[omDef,thick,const plot mark mid] table[x=x,y=daily]{figdata/derivatives/04-greeks-and-the-hedging-pnl/discrete.csv};
\addplot[omThm,thick,const plot mark mid] table[x=x,y=eight]{figdata/derivatives/04-greeks-and-the-hedging-pnl/discrete.csv};
\legend{couverture quotidienne (21 fois),couverture 8 fois par jour (168)}
\end{axis}
\end{tikzpicture}

\omcaption{Distribution du P\&L d'un \omterm{def:dv:greeks-and-the-hedging-pnl:cashgamma}{straddle} vendu à découvert à une volatilité de 20,
lorsque l'action réalise exactement 20, sur 20\,000 trajectoires simulées. Les deux
distributions sont centrées sur zéro~; leurs écarts types, 0.85 et 0.31, décroissent
comme l'inverse de la racine carrée du nombre de couvertures. Données~: le
tutoriel.}\label{fig:dv:greeks-and-the-hedging-pnl:discrete}
\end{omfigure}

\section{Volatilité de point mort}

\begin{definition}[Volatilité de point mort]\label{def:dv:greeks-and-the-hedging-pnl:be}
La \emph{volatilité de point mort}\index{volatilité de point mort} d'une option couverte est
la volatilité réalisée à laquelle, sur une période, son P\&L gamma couvre exactement son
thêta et son financement. Le mouvement de point mort est le mouvement quotidien correspondant du
sous-jacent~: $\tfrac12\Gamma\,\delta S^2=-\Theta\,\delta t$.
\end{definition}

Avec des taux nuls, l'identité gamma--thêta fait de la \omterm{def:dv:greeks-and-the-hedging-pnl:be}{volatilité de point mort}
la volatilité implicite à laquelle l'option est valorisée. Pour le \omterm{def:dv:greeks-and-the-hedging-pnl:cashgamma}{straddle} de l'exemple,
le thêta vaut $-0.0757$ par jour calendaire et le gamma 0.1381~: le mouvement de point mort
vaut $\sqrt{2\times0.0757/0.1381}=1.047$, soit l'écart-type d'une journée à une volatilité de 20
sur 365 jours. Les desks qui comptent le thêta en jours calendaires mais ne voient les mouvements
que les jours ouvrés doivent faire en sorte que le mouvement sur les jours ouvrés paie aussi le thêta du
week-end~: sur 252 jours ouvrés, le mouvement de point mort vaut 1.26. Le choix de l'horloge est
une décision de modélisation, traitée au chapitre 8.

\section{À quelle volatilité couvrir}

Un desk qui est certain que le titre réalisera 25 choisit tout de même la volatilité de son
delta. Deux choix sont naturels.

\begin{proposition}[Couverture à la volatilité implicite contre couverture à la volatilité réalisée]\label{prop:dv:greeks-and-the-hedging-pnl:which}
Supposons que l'option soit achetée à la volatilité implicite $\sigma$ et que le titre réalise une volatilité constante
$\sigma_{\mathrm{real}}>\sigma$. Couvrir en continu avec le delta à
$\sigma_{\mathrm{real}}$ verrouille le profit $V(\sigma_{\mathrm{real}})-V(\sigma)$ dès
l'origine, mais la valorisation au prix de marché à la volatilité implicite fluctue en cours de route. Couvrir avec le delta à
$\sigma$ rend le P\&L de chaque jour une fonction connue du rendement
$x$ de ce jour, à savoir $\tfrac12\Gamma S^2(x^2-\sigma^2\delta t)$, mais le total dépend de la
trajectoire au travers de la pondération par le gamma~; son espérance est la même.
\end{proposition}
\begin{proof}
Avec le delta à $\sigma_{\mathrm{real}}$, l'option couverte est un
\omterm{def:dv:no-arbitrage-and-the-fundamental-theorems:claim}{portefeuille de réplication} pour le modèle vrai~: sa valeur terminale est fixée et égale à sa valeur
à $\sigma_{\mathrm{real}}$ à l'origine. Avec le delta à $\sigma$, appliquer
la \cref{def:dv:greeks-and-the-hedging-pnl:hpnl} et la
\cref{prop:dv:greeks-and-the-hedging-pnl:expected}.
\end{proof}

En pratique, la volatilité réalisée n'est pas connue à l'avance, si bien que le second choix
est celui par défaut~: la couverture se fait à la volatilité implicite, le P\&L du livre s'explique chaque jour par le gamma,
le thêta et le \omterm{def:dv:greeks-and-the-hedging-pnl:greeks}{vega} chaque jour, et la vision de volatilité réalisée s'exprime dans la taille
de la position. La \Cref{fig:dv:greeks-and-the-hedging-pnl:which} montre les deux choix
sur le mois du desk, couvert quatre fois par jour~: même moyenne, un écart-type de
0.72 à la volatilité implicite et de 0.54 à la volatilité réalisée, ce dernier n'étant que du bruit de couverture discrète.

\begin{omfigure}
\begin{tikzpicture}
\begin{axis}[width=10.5cm,height=4.8cm,
  xlabel={P\&L à l'échéance par straddle},ylabel={densité},
  tick label style={font=\small},label style={font=\small},
  xmin=-5,xmax=3,ymin=0,ymax=1.0,grid=major,grid style={gray!25},
  legend columns=2,legend style={font=\footnotesize,at={(0.5,-0.3)},anchor=north,draw=none,fill=none,/tikz/every even column/.append style={column sep=8pt}},
  table/col sep=comma]
\addplot[omDef,thick,const plot mark mid] table[x=x,y=implied]{figdata/derivatives/04-greeks-and-the-hedging-pnl/which_vol.csv};
\addplot[omThm,thick,const plot mark mid] table[x=x,y=realised]{figdata/derivatives/04-greeks-and-the-hedging-pnl/which_vol.csv};
\draw[gray,dashed] (axis cs:-1.151,0) -- (axis cs:-1.151,1.0);
\legend{couverture à la volatilité implicite 20,couverture à la volatilité réalisée 25}
\end{axis}
\end{tikzpicture}

\omcaption{Un \omterm{def:dv:greeks-and-the-hedging-pnl:cashgamma}{straddle} vendu à 20 alors que le titre réalise 25, couvert quatre
fois par jour avec le delta à 20 ou à 25. Les deux distributions sont centrées sur la perte espérée
$-1.151$ (pointillés). Couvrir à la vraie volatilité supprime la dépendance à la trajectoire et
ne laisse que le bruit de couverture discrète. Données~: le
tutoriel.}\label{fig:dv:greeks-and-the-hedging-pnl:which}
\end{omfigure}

\begin{definition}[Bump-and-reprice]\label{def:dv:greeks-and-the-hedging-pnl:bump}
Le \emph{bump-and-reprice}\index{bump-and-reprice} calcule une grecque en revalorisant la
position avec un paramètre décalé et en prenant une différence finie~: le delta comme
$(V(S+h)-V(S-h))/2h$, le gamma comme $(V(S+h)-2V(S)+V(S-h))/h^2$, le \omterm{def:dv:greeks-and-the-hedging-pnl:greeks}{vega} avec un décalage de
d'un point de volatilité. La méthode fonctionne pour tout pricer, au prix de deux ou trois
valorisations par grecque.
\end{definition}

Le pas est un compromis (livre~4, chapitre~25)~: trop grand, la
différence mesure une corde et non une pente~; trop petit, l'annulation catastrophique
ne laisse que du bruit d'arrondi, surtout pour le gamma. Les desks utilisent un pas relatif
sur le spot d'environ 1~\% pour les rapports de risque, correspondant aux mouvements qui les intéressent, et un pas plus petit
lorsqu'une dérivée régulière est nécessaire. Pour les pricers de Monte Carlo, les revalorisations doivent
utiliser les mêmes nombres aléatoires, faute de quoi le bruit de deux estimations indépendantes noie la
différence (chapitre 23).

\section{Tutoriel~: un mois de couverture}\label{tut:dv:greeks-and-the-hedging-pnl:tut}

\begin{tutorial}
\textbf{Objectif.} Simuler le mois complet du desk~: vendre des \omterm{def:dv:greeks-and-the-hedging-pnl:cashgamma}{straddles}, couvrir selon
un calendrier, comparer le P\&L à la prédiction gamma--thêta, et mesurer le
bruit de couverture discrète et l'effet de la volatilité de couverture. \textbf{État final~:}
\cref{fig:dv:greeks-and-the-hedging-pnl:path,fig:dv:greeks-and-the-hedging-pnl:discrete,fig:dv:greeks-and-the-hedging-pnl:which}
et les chiffres du problème du week-end.
\begin{tutsteps}
\item \textbf{Le simulateur}, vectorisé sur 20\,000 trajectoires~: vendre à la volatilité implicite, couvrir
à une volatilité choisie, financer à $r$, régler le \omterm{def:dv:greeks-and-the-hedging-pnl:cashgamma}{straddle} à l'échéance.
\omcode{code/derivatives/04-greeks-and-the-hedging-pnl/python/dv_greeks.py}{46}{69}{Couverture en delta d'un straddle vendu sur de nombreuses trajectoires à la fois.}\label{lst:dv:greeks-and-the-hedging-pnl:sim}
\item \textbf{\omterm{def:dv:greeks-and-the-hedging-pnl:greeks}{Grecques} en unités de desk} par \omterm{def:dv:greeks-and-the-hedging-pnl:bump}{bump-and-reprice}, pour tout pricer~:
\omcode{code/firm/greeks/firm_greeks.py}{16}{30}{Grecques par bump-and-reprice centrées, rapportées par point, par jour et par point de base.}\label{lst:dv:greeks-and-the-hedging-pnl:bump}
\item \textbf{Exécuter} \texttt{dv\_greeks.problem()} et \texttt{fig\_greeks.py}~; comparer
les écarts-types simulés avec la \cref{prop:dv:greeks-and-the-hedging-pnl:dk}.
\end{tutsteps}
\textbf{Pour aller plus loin.} Ajouter un coût de bid--ask sur chaque transaction de couverture et trouver la
fréquence de couverture qui minimise le coût plus le bruit~; puis laisser la volatilité valoir 25 pendant la
première moitié du mois et 15 pendant la seconde, et comparer le P\&L d'un \omterm{def:dv:greeks-and-the-hedging-pnl:cashgamma}{straddle}
de strike 100 sur les trajectoires qui finissent à 100 et à 110.
\end{tutorial}

\section{Construction~: le calculateur de grecques}\label{bld:dv:greeks-and-the-hedging-pnl:greeks}

\begin{build}
\bfield{Rôle} Chaque position de la firme miniature rapporte les mêmes \omterm{def:dv:greeks-and-the-hedging-pnl:greeks}{grecques} dans
les mêmes unités, quel que soit le pricer qui la valorise~; le P\&L du lendemain est prédit à partir
de celles-ci, de sorte que la part non expliquée puisse être surveillée (chapitre 25). \bfield{Interface}
\texttt{bump\_greeks(pricer, params, h\_spot, h\_vol, h\_rate,
days)} avec \texttt{pricer(**params)}~; \texttt{desk\_units(greeks, spot)}~;
\texttt{predict\_pnl(greeks, d\_spot, d\_vol\_pts, days)}~;
\texttt{hedging\_pnl(cash\_gamma, ret, vol, dt)}~; \texttt{breakeven\_vol}.
\bfield{Règles} Différences centrées pour le delta, le gamma, le \omterm{def:dv:greeks-and-the-hedging-pnl:greeks}{vega} et le \omterm{def:dv:greeks-and-the-hedging-pnl:greeks}{rho}~; thêta en faisant avancer la
date de valorisation d'un jour calendaire~; unités~: actions, $\Gamma S^2/100$, par point de
volatilité, par jour, par point de base.
\bfield{Tests d'acceptation} \texttt{code/firm/greeks/tests/}~: le \omterm{def:dv:greeks-and-the-hedging-pnl:bump}{bump-and-reprice} égale
les \omterm{def:dv:greeks-and-the-hedging-pnl:greeks}{grecques} analytiques de \texttt{firm.bs} en unités de desk~; la prédiction à un jour est
à moins d'un centime d'une revalorisation complète pour un mouvement de 1~\%~; la \omterm{def:dv:greeks-and-the-hedging-pnl:be}{volatilité de point mort}
égale la volatilité implicite à taux nuls.
\bfield{Pour aller plus loin} Les mêmes \omterm{def:dv:greeks-and-the-hedging-pnl:greeks}{grecques} pour un portefeuille valorisé par le moteur de Monte Carlo du
chapitre 23, avec des nombres aléatoires communs~; la bibliothèque de pricing du chapitre 28 les fournit
pour tout instrument par un seul appel.
\end{build}

\begin{omsources}
\item P.~P.~Boyle et D.~Emanuel, «~Discretely adjusted option hedges~»,
\emph{Journal of Financial Economics} 8 (1980) 259--282.
\item E.~Derman et M.~Kamal, «~When you cannot hedge continuously~: the
corrections of Black--Scholes~», \emph{Risk} 12 (1999) 82--85.
\item R.~Ahmad et P.~Wilmott, «~Which free lunch would you like today, sir~?
Delta hedging, volatility arbitrage and optimal portfolios~», \emph{Wilmott}
(2005) 64--79.
\item H.~E.~Leland, «~Option pricing and replication with transactions costs~»,
\emph{Journal of Finance} 40 (1985) 1283--1301.
\end{omsources}

\section{Exercices}

\begin{exercise}[$\star$]\label{exo:dv:greeks-and-the-hedging-pnl:1}
Donner le delta, le gamma, le \omterm{def:dv:greeks-and-the-hedging-pnl:greeks}{vega} par point et le thêta par jour du \omterm{def:dv:greeks-and-the-hedging-pnl:cashgamma}{straddle} à la monnaie
d'un mois sur une action à 100 avec une volatilité de 20, taux nuls.
\end{exercise}

\begin{exercise}[$\star$]\label{exo:dv:greeks-and-the-hedging-pnl:2}
Vérifier l'identité gamma--thêta sur les valeurs de l'\cref{exo:dv:greeks-and-the-hedging-pnl:1}.
\end{exercise}

\begin{exercise}[$\star$]\label{exo:dv:greeks-and-the-hedging-pnl:3}
Une option longue couverte a un \omterm{def:dv:greeks-and-the-hedging-pnl:cashgamma}{cash gamma} de 5\,000 et est valorisée à une volatilité de 20.
Le titre se déplace de 2~\% aujourd'hui. Donner le P\&L de gamma du jour, le thêta et le résultat net
(taux nuls, un jour calendaire).
\end{exercise}

\begin{exercise}[$\star\star$]\label{exo:dv:greeks-and-the-hedging-pnl:4}
Pourquoi le \omterm{def:dv:greeks-and-the-hedging-pnl:hpnl}{P\&L de couverture} est-il aveugle au signe des mouvements~? Donner un chemin sur lequel
un \omterm{def:dv:greeks-and-the-hedging-pnl:cashgamma}{straddle} long perd de l'argent alors que le titre finit loin du prix d'exercice.
\end{exercise}

\begin{exercise}[$\star\star$]\label{exo:dv:greeks-and-the-hedging-pnl:5}
De combien le nombre de couvertures doit-il augmenter pour diviser par deux l'\omterm{def:dv:greeks-and-the-hedging-pnl:dhe}{erreur de couverture discrète}~?
Quel est le coût si chaque transaction de couverture paie une demi-fourchette~?
\end{exercise}

\begin{exercise}[$\star\star$]\label{exo:dv:greeks-and-the-hedging-pnl:6}
En utilisant la \cref{prop:dv:greeks-and-the-hedging-pnl:dk}, donner l'écart type de
l'erreur de couverture quotidienne du \omterm{def:dv:greeks-and-the-hedging-pnl:cashgamma}{straddle} du desk, et la comparer
au 0.865 simulé.
\end{exercise}

\begin{exercise}[$\star\star\star$]\label{exo:dv:greeks-and-the-hedging-pnl:7}
\emph{Codage.} Avec \texttt{hedge\_short\_straddle}, mesurer la moyenne et l'écart type du P\&L
lorsque le \omterm{def:dv:greeks-and-the-hedging-pnl:cashgamma}{straddle} est vendu à 20, que le titre réalise 25 et que la
couverture est à 25, couverte quotidiennement, avec la graine 1.
\end{exercise}

\begin{exercise}[$\star\star\star$]\label{exo:dv:greeks-and-the-hedging-pnl:8}
\emph{Repérer l'erreur.} «~Nos \omterm{def:dv:greeks-and-the-hedging-pnl:cashgamma}{straddles} courts ont rapporté sur 70~\% des jours ce mois-ci
et ont perdu au total. Le thêta jouait pour nous~; la perte doit venir d'une erreur de pricing.~»
\end{exercise}

\section{Problème~: le mois du straddle vendu}

\begin{problem}[{Problème du week-end --- vendre une volatilité qui s'est révélée trop bon marché}]\label{pb:dv:greeks-and-the-hedging-pnl:1}
Un desk vend 1\,000 \omterm{def:dv:greeks-and-the-hedging-pnl:cashgamma}{straddles} à la monnaie d'un mois (multiplicateur 100) sur une action à
100, à une volatilité implicite de 20, taux nuls et sans dividende, et couvre quotidiennement. Au cours du
mois, le titre réalise 25.

\medskip\noindent\textbf{Partie I --- La position.}
\begin{enumerate}
\item Donner la prime reçue, par \omterm{def:dv:greeks-and-the-hedging-pnl:cashgamma}{straddle} et en dollars.
\item Donner le delta de la position en actions et la couverture initiale.
\item Donner son \omterm{def:dv:greeks-and-the-hedging-pnl:cashgamma}{cash gamma} et son gamma pour un mouvement de 1~\%, par \omterm{def:dv:greeks-and-the-hedging-pnl:cashgamma}{straddle}.
\item Donner son \omterm{def:dv:greeks-and-the-hedging-pnl:greeks}{vega} en dollars par point de volatilité et son thêta par jour.
\item Donner le mouvement quotidien de point mort sur 365 jours puis sur 252 jours.
\end{enumerate}

\medskip\noindent\textbf{Partie II --- Le résultat attendu.}
\begin{enumerate}[resume]
\item Donner la valeur du \omterm{def:dv:greeks-and-the-hedging-pnl:cashgamma}{straddle} à une volatilité de 25.
\item Donner la perte attendue par \omterm{def:dv:greeks-and-the-hedging-pnl:cashgamma}{straddle} et en dollars.
\item La comparer à cinq fois le \omterm{def:dv:greeks-and-the-hedging-pnl:greeks}{vega}.
\item Pourquoi les deux résultats concordent-ils si bien~?
\item La réponse changerait-elle si le desk couvrait en continu~?
\end{enumerate}

\medskip\noindent\textbf{Partie III --- Le mois simulé.}
\begin{enumerate}[resume]
\item Donner la moyenne et l'écart type simulés du P\&L par \omterm{def:dv:greeks-and-the-hedging-pnl:cashgamma}{straddle} avec une couverture quotidienne.
\item Donner l'écart type avec huit couvertures par jour.
\item Donner la probabilité d'une perte avec une couverture quotidienne.
\item Donner la moyenne et l'écart type lorsque le desk couvre à 25 au lieu de 20.
\item Avec une volatilité réalisée égale à 20, quelles sont la moyenne et l'écart type~?
\end{enumerate}

\medskip\noindent\textbf{Partie IV --- Jugement.}
\begin{enumerate}[resume]
\item Que fallait-il que le desk croie pour vendre à 20~?
\item Comment dimensionner la position pour qu'un mauvais mois ne mette pas fin au desk~?
\item Quelle grecque du rapport aurait dû alerter le desk~?
\item Énoncer le \emph{résultat nommé}~: la perte attendue de la position à partir de l'identité
gamma--thêta et son écart type sous une couverture quotidienne, en dollars.
\item En une phrase~: que rapporte une option couverte en delta~?
\end{enumerate}
\end{problem}

\section{Questions d'entretien}

\begin{interviewq}[$\star$ \iqroles{trader}]\label{iq:dv:greeks-and-the-hedging-pnl:1}
Vous êtes long d'une option à la monnaie couverte en delta. Le sous-jacent ne bouge pas pendant
une semaine. Qu'est-il arrivé à votre P\&L, et pourquoi~?
\end{interviewq}

\begin{interviewq}[$\star$ \iqroles{trader, chercheur}]\label{iq:dv:greeks-and-the-hedging-pnl:2}
Énoncer la relation gamma--thêta et l'expliquer sans formules.
\end{interviewq}

\begin{interviewq}[$\star\star$ \iqroles{chercheur}]\label{iq:dv:greeks-and-the-hedging-pnl:3}
Établir le P\&L d'une option couverte en delta sur un jour en fonction de son gamma et des
volatilités réalisée et implicite.
\end{interviewq}

\begin{interviewq}[$\star\star$ \iqroles{trader}]\label{iq:dv:greeks-and-the-hedging-pnl:4}
Vous êtes certain que la volatilité réalisée sera supérieure à l'implicite. Faut-il couvrir en delta à l'implicite
ou à votre prévision~? Qu'est-ce qui change~?
\end{interviewq}

\begin{interviewq}[$\star\star$ \iqroles{développeur, risque}]\label{iq:dv:greeks-and-the-hedging-pnl:5}
Votre gamma en \omterm{def:dv:greeks-and-the-hedging-pnl:bump}{bump-and-reprice} est bruité pour un pricer de Monte Carlo et faux pour un digital
proche de l'échéance. Pourquoi, et que faites-vous~?
\end{interviewq}

\begin{interviewq}[$\star\star\star$ \iqroles{chercheur, trader}]\label{iq:dv:greeks-and-the-hedging-pnl:6}
Deux portefeuilles ont le même \omterm{def:dv:greeks-and-the-hedging-pnl:greeks}{vega} et la même différence entre volatilité réalisée et implicite sur un
mois. Pourquoi leurs P\&L peuvent-ils différer, et de combien~?
\end{interviewq}
